% Pursuing Stacks (cotes 134-2 to 134-8): archive cote 134-4, site page 27 (PDF page 28), typescript page 115.
% typescript. Draft compared: ps3.tex:1329-1392.
% AI-transcribed by Claude Opus 5.5 (claude-opus-5-5, adaptive thinking, effort high), Message Batches API, batch ps2 (msgbatch_01NVwr1Yuo8qJf5rykeGBXPK), 2026-11-12; not yet checked by eye.
\scanpage{134-4/27}
\mspage{115}

\inserted{to,} is a localization functor.

\xout{Here is a \uline{sufficient} condition on $R_h$ for an element $R_h$ of $\mathrm{Hom}_1$ to be}
\xout{in the image of $\mathrm{Hom}_2$ , namely to be the}
For a general $R_h \in \mathrm{Hom}_1(M)$ , if we consider the \unsure{comrresponding|corresponding} $W_h$
$( = r_{21}(R_h))$ in $\mathrm{Hom}_2(M)$, it is clear that the canonical functor $M \to M_h$
of A factors into
\[ M \to W_h^{-1}M \to M_h \ , \]
and $R_h$ ``is in $\mathrm{Hom}_2(M)$'' i.e. $R_h = r_{12}(W_h)$, iff the second functor
\[ W_h^{-1}M \to M_h \]
(which is anyhow bijective on objects and surjective on arrows) is an iso-%
morphism, or, equivalently, faithful. Here is a rather evident \deleted{\unsure{and}}
\uline{sufficient} condition on $R_h$ for this to be so,\hand{$^*$} namely\hand{:}
\margin{* 3.4. See \S 52 for a \unsure{proof} of this rash statement!}
\xout{\unsure{(C}} C$_{12}$) If $f,\unsure{g} : X \to Y$ are homotopic, there exists a homotopism
$X' \to X$ , \deleted{and} two sections $s_o,s_1$ of $X'$ over $X$, and a map $h \colon X' \to Y$ ,
such that
\[ f_o = hs_o \ , \ f_1 = hs_1 \ . \]
\uline{Remark.}Intuitively, we are thinking of course of $X'$ as a product $X\mathrm{x}I$ , where
$\underline{I} = (I,e_o,e_1)$ is a \inserted{weak} homotopy interval, and $s_o,s_1$ are defined in terms of
$e_o,e_1$ . In Quillen's somewhat different set-up, $X'$ is referred to as a
``cylinder object for $X$'' , suitable for defining the ``left homotopy rela-%
tion'' associated to a given $W_h$ . The condition C$_{12}$ is not autodual,
we could state a dual sufficient condition in terms of a ``path object for
$Y$ '', namely a homotopism $Y \to Y'$ endowed with two retractions $t_o,t_1$ upon
$Y$ - but we won't have any use for this \xout{xxx}\inserted[typed]{in} the present set up, which
(as far as the main emphasis is concerned) is by no means autodual, as
is Quillen's.

\margin{\unsure{$\cdot$}} The condition C$_{12}$ above can be viewed equally as a condition on
a $W_h \in \mathrm{Hom}_2(M)$.

We may interprete the set $\mathrm{Hom}_2(M)$ of homotopism notions in $M$ as the
set of all quotient categories $M_h$ of $M$, having same \xout{xxxxxx} objects as $M$,
and such that moreover the canonical functor $M \to M_h$ be a localizing func-%
tor . As in A), the relevant ``\uline{accessory assumption}'' on $W_h$ (a particular
case indeed of the corresponding one for $R_h$) is that this functor commute
to products. I dont see any \xout{handy} simple computational way though to \xout{expr}
express this condition directly in terms of $W_h$ , as previously in terms of
$R_h$. I would only like to notice here a consequence of this assumption
(I doubt it is equivalent to it), namely that \uline{the cartesian product of
two homotopisms is again a homotopism} - which implies, for instance,

\margin{27}
